%--------------------------------------------------------------------
\subsection{Current Operator and Nonlocal Pseudopotentials}
\label{sec:current-nlpp}
%--------------------------------------------------------------------

For a periodic system in the homogeneous-field limit, the optical
response is naturally formulated in terms of the macroscopic
charge-current density.  We define
\begin{equation}
  J_\alpha(t)
  =
  \frac{1}{V_{\rm cell}}
  \sum_{\mathbf k,\sigma}
  w_{\mathbf k}\,
  {\rm Tr}\!
  \left[
    \mathbf P_{\mathbf k}^{\sigma}(t)
    \hat I_{\mathbf k,\alpha}^{\sigma}(t)
  \right],
  \label{eq:macroscopic_current_density}
\end{equation}
where $V_{\rm cell}$ is the volume of the periodic simulation cell,
which may be a primitive cell or a supercell, $w_{\mathbf k}$ are the
normalized Brillouin-zone weights satisfying
$\sum_{\mathbf k}w_{\mathbf k}=1$, $\mathbf P_{\mathbf k}^{\sigma}$ is
the single-particle density matrix for spin channel $\sigma$, and
\begin{equation}
  \hat I_\alpha(t)=q_e\hat v_\alpha(t)
  \label{eq:charge_current_operator_definition}
\end{equation}
is the charge-current operator.  The trace in
Eq.~\eqref{eq:macroscopic_current_density} is taken over the propagated
single-particle subspace at fixed $\mathbf k$ and spin.

The velocity operator is most cleanly defined before introducing the
cell-periodic Bloch representation.  For the velocity-gauge Hamiltonian
of Eq.~\eqref{eq:velocity_gauge_hamiltonian},
\begin{equation}
  \hat v_\alpha(t)
  =
  \frac{\imath}{\hbar}
  \left[
    \hat H_V(t),\hat r_\alpha
  \right].
  \label{eq:velocity_commutator_definition}
\end{equation}
Since the local potential commutes with the position operator, the
Hamiltonian containing a local potential and a nonlocal pseudopotential
gives
\begin{equation}
  \hat v_\alpha(t)
  =
  \frac{1}{m_e}
  \left[
    \hat p_\alpha
    -
    \frac{q_e}{c}A_\alpha(t)
  \right]
  +
  \frac{\imath}{\hbar}
  \left[
    \hat V_{\rm NL}^{V}(t),\hat r_\alpha
  \right].
  \label{eq:velocity_operator_full}
\end{equation}
Equivalently,
\begin{equation}
  \frac{\imath}{\hbar}
  \left[
    \hat V_{\rm NL}^{V}(t),\hat r_\alpha
  \right]
  =
  \frac{1}{\imath\hbar}
  \left[
    \hat r_\alpha,\hat V_{\rm NL}^{V}(t)
  \right],
  \label{eq:nlpp_velocity_commutator_equivalence}
\end{equation}
which fixes the sign convention for the nonlocal contribution.  Thus the
formal charge-current operator is
\begin{equation}
  \hat I_\alpha(t)
  =
  q_e
  \left\{
  \frac{1}{m_e}
  \left[
    \hat p_\alpha
    -
    \frac{q_e}{c}A_\alpha(t)
  \right]
  +
  \frac{\imath}{\hbar}
  \left[
    \hat V_{\rm NL}^{V}(t),\hat r_\alpha
  \right]
  \right\}.
  \label{eq:formal_charge_current_operator}
\end{equation}
The explicit $A_\alpha(t)$ dependence of the kinetic term gives the
usual diamagnetic contribution,
\begin{equation}
  \hat I_{\alpha}^{\rm dia}(t)
  =
  -
  \frac{q_e^2}{m_e c}
  A_\alpha(t)\hat 1,
  \label{eq:diamagnetic_current_operator}
\end{equation}
while the nonlocal term is also $A$ dependent through the
gauge-transformed nonlocal pseudopotential.

The same nonlocal commutator appears when the velocity-gauge Hamiltonian
is expanded to first order in the homogeneous vector potential.  Using
Eq.~\eqref{eq:gauge_transformed_nlpp},
\begin{equation}
  \hat V_{\rm NL}^{V}(t)
  =
  \hat U(t)\hat V_{\rm NL}\hat U^\dagger(t),
  \qquad
  \hat U(t)
  =
  \exp\!\left[
    \frac{\imath q_e}{\hbar c}
    \mathbf r\cdot\mathbf A(t)
  \right],
  \label{eq:nlpp_gauge_transform_repeated}
\end{equation}
one obtains
\begin{equation}
  \hat V_{\rm NL}^{V}(t)
  =
  \hat V_{\rm NL}
  +
  \frac{\imath q_e}{\hbar c}
  \sum_\alpha
  A_\alpha(t)
  \left[
    \hat r_\alpha,\hat V_{\rm NL}
  \right]
  +
  {\cal O}(A^2).
  \label{eq:nlpp_gauge_transform_linear}
\end{equation}
The first-order nonlocal perturbation is therefore
\begin{equation}
  \delta\hat V_{\rm NL}^{(1)}(t)
  =
  -
  \frac{1}{c}
  \sum_\alpha
  A_\alpha(t)\,
  \hat I_{\alpha,{\rm NL}}^{(0)},
  \qquad
  \hat I_{\alpha,{\rm NL}}^{(0)}
  =
  \frac{q_e}{\imath\hbar}
  \left[
    \hat r_\alpha,\hat V_{\rm NL}
  \right].
  \label{eq:nlpp_linear_perturbation_current}
\end{equation}
Combining this term with the kinetic part of the minimal-coupling
Hamiltonian gives the linear velocity-gauge perturbation
\begin{equation}
  \delta\hat H^{(1)}(t)
  =
  -
  \frac{1}{c}
  \sum_\alpha
  A_\alpha(t)\,
  \hat I_{\alpha}^{(0)},
  \label{eq:linear_velocity_gauge_perturbation}
\end{equation}
where
\begin{equation}
  \hat I_{\alpha}^{(0)}
  =
  q_e
  \left[
    \frac{\hat p_\alpha}{m_e}
    +
    \frac{\imath}{\hbar}
    \left[
      \hat V_{\rm NL},\hat r_\alpha
    \right]
  \right]
  =
  \frac{q_e}{m_e}\hat p_\alpha
  +
  \frac{q_e}{\imath\hbar}
  \left[
    \hat r_\alpha,\hat V_{\rm NL}
  \right].
  \label{eq:zero_field_charge_current_operator}
\end{equation}
Equations~\eqref{eq:formal_charge_current_operator}--\eqref{eq:zero_field_charge_current_operator}
show that the nonlocal-pseudopotential commutator is not only a
correction to the measured current; it is also the nonlocal part of the
linear velocity-gauge perturbation required for gauge consistency when
nonlocal pseudopotentials are present.\cite{Luber-rt-tddft}  More
generally, a nonlocal one-particle operator present in the Hamiltonian
contributes to the velocity through the corresponding commutator with the
position operator.

For a separable norm-conserving pseudopotential,\cite{KleinmanBylander}
\begin{equation}
  \hat V_{\rm NL}
  =
  \sum_{I,lm}
  |\beta_{I,lm}\rangle
  D_{I,lm}
  \langle\beta_{I,lm}|,
  \label{eq:separable_nlpp}
\end{equation}
the coordinate-space matrix element of the commutator can be written as
\begin{equation}
  \langle\mathbf r|
  \left[
    \hat r_\alpha,\hat V_{\rm NL}
  \right]
  |\mathbf r'\rangle
  =
  (r_\alpha-r'_\alpha)
  V_{\rm NL}(\mathbf r,\mathbf r').
  \label{eq:nlpp_commutator_coordinate}
\end{equation}
Equation~\eqref{eq:nlpp_commutator_coordinate} makes explicit that this
term vanishes for a strictly local potential but is finite for a
nonlocal pseudopotential.

The corresponding cell-periodic Bloch representation is obtained only
after the full-space current operator has been defined.  For
$\Psi_{n\mathbf k}=e^{\imath\mathbf k\cdot\mathbf r}
u_{n\mathbf k}$, the canonical momentum becomes
$-\imath\hbar\nabla+\hbar\mathbf k$ as in
Eq.~\eqref{eq:bloch_momentum_operator}.  The nonlocal operator at fixed
$\mathbf k$ is
\begin{equation}
  \hat V_{{\rm NL},\mathbf k}
  =
  e^{-\imath\mathbf k\cdot\mathbf r}
  \hat V_{\rm NL}
  e^{\imath\mathbf k\cdot\mathbf r}.
  \label{eq:nlpp_bloch_transform}
\end{equation}
At zero external vector potential, the charge-current matrix in a
cell-periodic basis is therefore
\begin{align}
  I_{ab,\mathbf k,\alpha}^{(0),\sigma}
  &=
  \frac{q_e}{m_e}
  \left\langle
    u_{a\mathbf k}^{\sigma}
    \left|
    -\imath\hbar\nabla_\alpha+\hbar k_\alpha
    \right|
    u_{b\mathbf k}^{\sigma}
  \right\rangle_{\rm cell}
  \nonumber\\
  &\quad
  +
  \frac{q_e}{\imath\hbar}
  \left\langle
    u_{a\mathbf k}^{\sigma}
    \left|
    \left[
      \hat r_\alpha,\hat V_{{\rm NL},\mathbf k}^{\sigma}
    \right]
    \right|
    u_{b\mathbf k}^{\sigma}
  \right\rangle_{\rm cell}.
  \label{eq:bloch_current_matrix}
\end{align}
The first term contains both the gradient and Bloch
$\hbar k_\alpha$ contributions, while the second term is the
nonlocal-pseudopotential contribution.  In the presence of a finite
homogeneous vector potential, Eq.~\eqref{eq:formal_charge_current_operator}
gives the corresponding $A$-dependent current operator, including the
diamagnetic kinetic term and the gauge-transformed nonlocal operator.

Equations~\eqref{eq:macroscopic_current_density} and
\eqref{eq:bloch_current_matrix} define the formal current observable used
in the periodic RT-TDDFT theory.  The implementation-specific
construction of the current matrices, the particular impulsive
realization of the velocity-gauge perturbation, and the conversion of raw
current traces to optical response functions are described separately in
Secs.~\ref{sec:implementation} and~\ref{sec:optical-response}.
